\documentclass[11pt]{article}
\usepackage[margin=1in]{geometry}
\usepackage{amsmath,amssymb,enumitem}
\usepackage[T1]{fontenc}
\usepackage[utf8]{inputenc}
\usepackage[colorlinks=true,linkcolor=blue,urlcolor=blue,citecolor=blue]{hyperref}

\title{Both Closing Problems in \texttt{arXiv:1301.5799}: Literature Status}
\author{\texttt{fa\_banach\_001}, \texttt{agent\_lane\_12}}
\date{13 August 2026}

\begin{document}
\maketitle

\section*{Status}

\textbf{Literature already answered.}
The two closing questions in Cuth's \emph{Noncommutative Valdivia compacta}
are answered in separate later papers.  The first answer identifies the
requested function-space property, with a countable-compactness condition in
the merely dense formulation; the second gives a four-way equivalence
stronger than the three requested implications.

\section*{Source: Problem 1}

On PDF page 13 of arXiv:1301.5799, Problem 1 asks for a topological property
$(T)$ of $(C(K),\tau_p(D))$ which characterizes those dense (respectively,
dense and countably closed) sets $D\subset K$ induced by a retractional
skeleton.

Cuth--Kalenda, \emph{Monotone retractability and retractional skeletons},
arXiv:1403.4480, explicitly says that its Theorem 1.4 answers this problem.
That theorem (PDF page 2; proof on page 11) states that, for compact $K$ and
dense $D\subset K$, the following are equivalent:
\begin{enumerate}[label=(\roman*)]
  \item $D$ is induced by a retractional skeleton in $K$;
  \item $D$ is countably compact and $(C(K),\tau_p(D))$ is monotonically
        Sokolov.
\end{enumerate}
Thus the sought function-space property is \emph{monotone Sokolovness}.  In
the source's dense-and-countably-closed branch, $D$ is automatically
countably compact.  For a merely dense $D$, the exact characterization keeps
countable compactness as an additional condition on $D$.

\newpage
\section*{Source: Question 1}

The same source page asks whether, for compact $K$, the following implications
hold among:
\begin{enumerate}[label=(\roman*)]
  \item $C(K)$ has a $1$-projectional skeleton;
  \item a convex symmetric set is induced by a retractional skeleton in
        $(B_{C(K)^*},w^*)$;
  \item a convex symmetric set is induced by a retractional skeleton in
        $P(K)$.
\end{enumerate}
Specifically, it asks for (iii)$\Rightarrow$(ii), (ii)$\Rightarrow$(i), and
(iii)$\Rightarrow$(i).

Cuth, \emph{Simultaneous projectional skeletons}, arXiv:1305.1438, states
that Theorem 4.1 answers this question.  The theorem (PDF page 6) proves the
equivalence of:
\begin{enumerate}[label=(\roman*)]
  \item $C(K)$ has a $1$-projectional skeleton;
  \item there is a convex symmetric set induced by a retractional skeleton in
        $(B_{C(K)^*},w^*)$;
  \item there is a convex set induced by such a skeleton in the same dual ball;
  \item there is a convex set induced by a retractional skeleton in $P(K)$.
\end{enumerate}
Condition (iv) is weaker than the source's symmetric $P(K)$ condition.
Therefore all three implications asked in Question 1 have affirmative
answers.

\section*{Scope}

This is a literature-status identification, not a new proof.  Both answering
papers explicitly cite and answer the numbered problem or question in
arXiv:1301.5799.  The dense-only formulation above is reported with the exact
countable-compactness qualification of Theorem 1.4.

\begin{thebibliography}{9}
\bibitem{Cuth2013}
M. Cuth,
\emph{Noncommutative Valdivia compacta},
Comment. Math. Univ. Carolin. 55 (2014), no. 1, 53--72;
arXiv:1301.5799.

\bibitem{Cuth2014}
M. Cuth,
\emph{Simultaneous projectional skeletons},
J. Math. Anal. Appl. 411 (2014), 19--29;
arXiv:1305.1438.

\bibitem{CK2015}
M. Cuth and O. F. K. Kalenda,
\emph{Monotone retractability and retractional skeletons},
J. Math. Anal. Appl. 423 (2015), 18--31;
arXiv:1403.4480.
\end{thebibliography}

\end{document}
